\chapter{Arithmetization of Syntax}

\section{G\"odel Numbering}

\begin{defn}
    For the formal system of Peano Arithmetic, the \textsl{syntax} is the coproduct of the inductive types for terms and formulas:
    \[
        \fun{Syntax}\defeq\fun{Term}+\fun{Form}.
    \]
    Its \textsl{G\"odel numbering} is a function assigning a unique natural number to each syntactic object:
    \begin{align*}
        \fun{gn}\colon\fun{Syntax}&\to\N\\
        s&\mapsto\gn{s}.
    \end{align*}
\end{defn}

For this coding to be sound, the numbering function must be injective.

To construct the G\"odel numbering constructively, we rely on the Fundamental Theorem of Arithmetic. By mapping each syntactic constructor to a unique power of $2$, and encoding its sub-components into subsequent prime factors ($3$ and $5$), we guarantee that every expression maps to a unique natural number.

We construct the mapping in two stages: first by defining structural recursors for terms and formulas, and then by uniting them via the coproduct eliminator for the syntax type.

\begin{defn}
    We define the functions
    \[
        \fun{gn}_{\fun{Term}}\colon\fun{Term}\to\N
        \quad\text{and}\quad
        \fun{gn}_{\fun{Form}}\colon\fun{Form}\to\N
    \]
    by structural recursion. 
    
    For terms,
    \begin{align*}
        \fun{gn}_{\fun{Term}}(\fun{var}(x))
            &\defeq 2^1\cdot3^x\\
        \fun{gn}_{\fun{Term}}(0)
            &\defeq 2^2\\
        \fun{gn}_{\fun{Term}}(s(t))
            &\defeq 2^3\cdot3^{\fun{gn}_{\fun{Term}}(t)}\\
        \fun{gn}_{\fun{Term}}(t_1+t_2)
            &\defeq
                2^4
                \cdot
                3^{\fun{gn}_{\fun{Term}}(t_1)}
                \cdot
                5^{\fun{gn}_{\fun{Term}}(t_2)}\\
        \fun{gn}_{\fun{Term}}(t_1\cdot t_2)
            &\defeq
                2^5
                \cdot
                3^{\fun{gn}_{\fun{Term}}(t_1)}
                \cdot
                5^{\fun{gn}_{\fun{Term}}(t_2)}
    \intertext{For formulas,}
        \fun{gn}_{\fun{Form}}(t_1\doteq t_2)
            &\defeq
                2^6
                \cdot
                3^{\fun{gn}_{\fun{Term}}(t_1)}
                \cdot 5^{\fun{gn}_{\fun{Term}}(t_2)} \\
        \fun{gn}_{\fun{Form}}(\neg\varphi)
            &\defeq
                2^7
                \cdot
                3^{\fun{gn}_{\fun{Form}}(\varphi)}\\
        \fun{gn}_{\fun{Form}}(\varphi\imp\psi)
            &\defeq
                2^8
                \cdot
                3^{\fun{gn}_{\fun{Form}}(\varphi)}
                \cdot
                5^{\fun{gn}_{\fun{Form}}(\psi)}\\
        \fun{gn}_{\fun{Form}}(\forall x\,\varphi)
            &\defeq
                2^9
                \cdot
                3^x
                \cdot 5^{\fun{gn}_{\fun{Form}}(\varphi)}.
    \end{align*}
    Since the type $\fun{Syntax}$ is defined as the coproduct $\fun{Term}+\fun{Form}$, we construct the global G\"odel numbering function $\fun{gn}\colon\fun{Syntax}\to\N$ by applying the non-dependent eliminator for coproducts. For any $c\colon\fun{Syntax}$:
    \[
        \fun{gn}(c)\defeq 
        \begin{cases}
            \fun{gn}_{\fun{Term}}(t)
                &\text{if }c\equiv\fun{inl}(t),\\
            \fun{gn}_{\fun{Form}}(\varphi)
                &\text{if }c\equiv\fun{inr}(\varphi).
        \end{cases}
    \]
\end{defn}

\begin{thm}
    The function\/ $\fun{gn}$ is injective. We express this by inhabiting the following type:\/
    \[
        \fun{isInj}(\fun{gn})\defeq
        \prod_{s_1,s_2\colon\fun{Syntax}} 
        (\gn{s_1}\eq{\N}\gn{s_2})
        \to (s_1 \eq{\fun{Syntax}}s_2).
    \]
\end{thm}

Furthermore, because the syntax is generated inductively and finitely, we can decide whether any given natural number is the G\"odel number of a valid syntactic object. We express this by stating that the image of the function $\gn{\dummy}$ is decidable over the natural numbers.

\begin{thm}
    The image of the function\/ $\fun{gn}$ is decidable over the natural numbers. Formally, we express this by inhabiting the type:\/
    \[
        \prod_{n\colon\N}\;
        \Bigl(
        \sum_{s\colon\fun{Syntax}}
        (\gn{s}\eq{\N} n)
        \Bigr) 
        + 
        \Bigl(
        \Bigl(
        \sum_{s\colon\fun{Syntax}}
        (\gn{s}\eq{\N} n)
        \Bigr)
        \to\fun0
        \Bigr).
    \]
\end{thm}

With the syntax encoded as natural numbers, we can arithmetize the derivation family. A derivation is encoded as a sequence of numbers, which itself is assigned a G\"odel number. We define the proof relation as a type family over pairs of natural numbers.

\begin{defn}
    The \textsl{proof predicate} is a type family taking a non-dependent pair of natural numbers, where the first represents the G\"odel number of a candidate derivation and the second represents the G\"odel number of a sentence:
    \[
        \fun{Proof}\colon\N\times\N\to\univ.
    \]
\end{defn}

\begin{thm}
    For every\/ $n,m\colon\N$, the type\/ $\fun{Proof}(n,m)$ is a mere proposition.
\end{thm}

\begin{thm}
    The proof relation is decidable over all non-dependent pairs. We express this by inhabiting the type:\/
    \[
        \prod_{n,m\colon\N}
        \fun{Proof}(n,m)+(\fun{Proof}(n,m)\to\fun0).
    \]
\end{thm}

\begin{defn}
    The \textsl{provability type} for a sentence with G\"odel number $m$ is defined as the dependent sum over the proof predicate:
    \[
        \fun{Provable}(m)\defeq\sum_{n\colon\N}\fun{Proof}(n,m).
    \]
    A valid proof witness inhabiting this type is a dependent pair:
    \[
        \dep{n,p}\colon\sum_{n\colon\N}\fun{Proof}(n,m),
    \]
    where $n$ is the G\"odel number of the derivation and $p$ is the computational evidence that $n$ genuinely proves $m$.
\end{defn}